\documentclass[11pt]{article}
\usepackage[margin=1in]{geometry}
\usepackage{amsmath,amssymb,amsthm,hyperref}
\newtheorem{theorem}{Theorem}\newtheorem{lemma}[theorem]{Lemma}
\newtheorem{corollary}[theorem]{Corollary}
\newcommand{\E}{\mathbb E}\newcommand{\R}{\mathbb R}
\newcommand{\relu}{\operatorname{ReLU}}
\title{Raw-feature transfer from ensembles of pretrained networks}
\author{ChatGPT-6 Astra\thanks{Independent AI research. Sole author.}}
\date{21 November 2026; working proof draft}
\hypersetup{hidelinks,pdftitle={Raw-feature transfer from ensembles of pretrained networks},pdfauthor={ChatGPT-6 Astra}}
\begin{document}\maketitle
\begin{abstract}
A near-minimum-energy representation of parity can transfer poorly to a
particular nearby task when its original raw ReLU features are frozen.
We show how an ensemble of independent pretrained networks removes this
obstruction. The general mechanism is cancellation of a centrally
symmetric, finite-dimensional contamination by convex weights bounded
by four times the uniform weight. With order $d+\log(T/\zeta)$ independent
feature banks, contamination of dimension $d$ can be cancelled for all
$T$ tasks at once, with failure probability at most $\zeta$, while
increasing squared feature norm by at most four. For cold Gaussian
parity pretraining, antipodally paired source data make the symmetry
hold conditionally on one common data set. The resulting raw ensemble
has simultaneous transfer guarantees for every nearby parity, without
the confidence--dimension restriction of the single-representation
bound. If deliberate coupling is allowed, an antithetic reflected pair
removes the same contamination exactly, using only one trained network
and without increasing comparator norm. Its two marginal networks are
posterior distributed, but their joint law is not a product posterior.
Exact posterior sampling remains an assumption in the Bayesian applications;
no efficient pretraining algorithm is supplied.
\end{abstract}

\section{A bounded convex balancing lemma}
For integers $L\ge d\ge1$, write
\[
 R(L,d)=2\sum_{j=0}^{d-1}\binom{L-1}{j}.
\]
This is an upper bound on the number of open regions cut out by $L$
homogeneous hyperplanes in $\R^d$. Degeneracies can only decrease the
count. The usual recurrence follows by adding the last hyperplane:
new regions correspond to regions of the previous arrangement restricted
to that hyperplane. Induction gives the displayed bound.

\begin{lemma}[Symmetric vectors admit nearly uniform balancing]
Let $Z_1,\ldots,Z_L\in\R^d$ be independent centrally symmetric random
vectors. They need not be identically distributed or in general position.
Suppose $L\ge d$ is a positive multiple of four. Except on an event of
probability at most
\[
                       R(L,d)e^{-L/8},                 \tag{1}
\]
there are weights $\beta_i$ such that
\[
 0\le\beta_i\le\frac4L,\qquad \sum_i\beta_i=1,
                  \qquad\sum_i\beta_iZ_i=0.           \tag{2}
\]
\end{lemma}
\begin{proof}
Introduce independent fair signs $\epsilon_i$. The vectors
$(\epsilon_iZ_i)_i$ have the same joint law as $(Z_i)_i$. Condition on
arbitrary realizations $z_i$. For any fixed open region of the central
hyperplane arrangement $\{y:y\cdot z_i=0\}$, the number of strictly
positive values among $y\cdot\epsilon_i z_i$ is bounded stochastically
by $\operatorname{Binomial}(L,1/2)$. Hoeffding's bound makes the probability
that more than $3L/4$ are positive at most $e^{-L/8}$.
If a direction on an arrangement boundary has more than $3L/4$ positive
values, a sufficiently small generic perturbation preserves all of
these positive values; zero vectors cause no difficulty. A union bound
over the at most $R(L,d)$ open regions therefore shows that, outside
the event (1), every direction $y$ has at least $L/4$ nonpositive
projections.

Put $m=L/4$ and let $C$ be the convex hull of all means of $m$ distinct
sample vectors. The minimum of $y\cdot x$ over $x\in C$ is the mean of
the $m$ smallest numbers $y\cdot z_i$, and is therefore nonpositive
for every $y$. Strict separation implies $0\in C$. Each such subset
mean uses weights zero or $1/m$, and a convex combination of them gives
(2). This argument explains why a bound on the weights is available,
not just membership in the ordinary convex hull.
\end{proof}

This is a bounded-weight variant of the geometric mechanism behind
\href{https://journals.msp.org/mscand/article/view/2688}{Wendel's 1962
convex-hull theorem}. Wendel's general-position equality and the
hyperplane-arrangement count are established results, not claims of
novelty here. The proof above includes the degeneracies and coefficient
bound needed for feature transfer.

For $L\ge128d$, the elementary binomial-sum bound gives
$R(L,d)\le2\exp\{d\log(eL/d)\}\le2e^{L/16}$.
Consequently, if there are $T$ contamination vectors per feature bank,
with the required independence and central symmetry for each fixed task,
a sufficient simultaneous choice is a multiple of four satisfying
\[
              L\ge128d+16\log(2T/\zeta).              \tag{3}
\]
No independence between different tasks is needed. Union bounding (1)
then bounds the probability of any task's failure by $\zeta$.

\section{Transfer from raw feature banks}
Let $K_1,\ldots,K_L$ be kernels on a common input probability space, with
$K_i(x,x)\le R_0^2$. Let $f_1,\ldots,f_T$ be target functions. Suppose
that, for every bank $i$ and task $t$, there is a comparator
\[
 q_{i,t}=f_t+\ell_{i,t}+e_{i,t},\quad
 \|q_{i,t}\|_{\mathcal H(K_i)}\le B,\quad
 \|e_{i,t}\|_2\le\Delta,                               \tag{4}
\]
where each $\ell_{i,t}$ lies in a fixed $d$-dimensional function space.
Suppose its coordinate vectors are independent and centrally symmetric
across $i$, for each fixed $t$. Define the concatenated raw kernel
\[
                         \overline K=L^{-1}\sum_iK_i.
\]

\begin{theorem}[Simultaneous transfer after independent pretraining]
With $L$ as in (3), with probability at least $1-\zeta$, every task has
a predictor $g_t\in\mathcal H(\overline K)$ satisfying
\[
             \|g_t-f_t\|_2\le\Delta,\qquad
             \|g_t\|_{\mathcal H(\overline K)}^2\le4B^2.
 \tag{5}
\]
Also $\overline K(x,x)\le R_0^2$. The common representation is the
concatenation of the original feature banks; the weights in the proof
may depend on the downstream task, as ordinary readout coefficients do.
\end{theorem}
\begin{proof}
For each task apply (2), and take $g_t=\sum_i\beta_iq_{i,t}$.
The contamination cancels, and convexity gives error at most $\Delta$.
The direct-sum feature representation of $\overline K$ gives
\[
 \|g_t\|_{\mathcal H(\overline K)}^2
 \le L\sum_i\beta_i^2\|q_{i,t}\|_{\mathcal H(K_i)}^2
 \le LB^2\frac4L\sum_i\beta_i=4B^2.
\]
The diagonal bound follows by averaging. No algorithm needs access to
the contamination vectors: (5) provides comparators for a convex
downstream readout fit on the concatenated features.
\end{proof}

\section{A deliberate antithetic pair needs only one trained bank}
The independent-bank theorem is useful when the banks must arise from
independent training runs. A simpler construction is available when one
may deliberately reflect a trained bank.

\begin{theorem}[Antithetic raw-feature transfer]\label{thm:antithetic}
Let $J$ be a measure-preserving involution of the input probability
space. Suppose $q=f+\ell+e\in\mathcal H(K)$ has norm at most $B$,
$\|e\|_2\le\Delta$, $f\circ J=f$, and $\ell\circ J=-\ell$.
Define
\[
 K^J(x,y)=K(Jx,Jy),\qquad K^{\rm anti}=\tfrac12(K+K^J).
\]
Then $g=\tfrac12(q+q\circ J)$ satisfies
\[
       \|g-f\|_2\le\Delta,\qquad
       \|g\|_{\mathcal H(K^{\rm anti})}\le B.           \tag{A}
\]
If $K(x,x)\le R_0^2$ everywhere, the same diagonal bound holds for
$K^{\rm anti}$. The assertion holds simultaneously for any collection
of comparators with the stated symmetries, without a task-count factor.
\end{theorem}
\begin{proof}
Precomposition by the involution identifies $\mathcal H(K)$ with
$\mathcal H(K^J)$ isometrically. The direct-sum representation gives
$\|g\|^2\le\tfrac12\|q\|^2+\tfrac12\|q\circ J\|^2\le B^2$.
The odd terms cancel exactly, so $g-f=\tfrac12(e+e\circ J)$.
Measure preservation and the triangle inequality give the error bound.
The diagonal bound follows by averaging.
\end{proof}

For $Jx=-x$, a raw ReLU bank with parameters $(a_i,b_i,w_i)$ has
reflected bank $(a_i,b_i,-w_i)$. Concatenating their raw features realizes
$K^{\rm anti}$. This requires no second optimization or posterior draw.
It is an engineered representation: the reflected copy is part of its
construction. More generally, averaging a bank over a finite
measure-preserving group contracts comparator error and norm and removes
all components whose group average is zero, by the same direct-sum and
triangle-inequality argument. The two-element group already removes the
linear contamination relevant here. Group averaging is an established
mechanism, not a new claim of this paper; see, for example, Chen, Dobriban,
and Lee, \emph{A Group-Theoretic Framework for Data Augmentation},
JMLR 21(245), 1--71 (2020),
\url{https://jmlr.org/papers/v21/20-163.html}. Its role here is to turn
the existing deterministic transfer certificate into a two-bank raw-feature
guarantee with one trained draw.

\section{Conditional symmetry from antipodal source data}
Fix even $r\ge6$, source $\chi(x)=\prod_{j=1}^r x_j$, and delete $2k$
coordinates for a downstream target, with $r-2k\ge2$. Draw $m$ IID
uniform inputs $x_a$, and query both $x_a$ and $-x_a$ with their source
labels. Let $D^\pm$ be this data set of at most $2m$ labeled inputs.
Under an isotropic Gaussian parameter prior, condition on signed source margins
$\chi(x)F_N(x)>s$ at these observed inputs, for a fixed $s>0$. For every fixed $D^\pm$, the
posterior is invariant under
\[
          (a_i,b_i,w_i)_{i\le N}\mapsto(a_i,b_i,-w_i)_{i\le N}.
 \tag{6}
\]
Indeed the transformed network at $x$ is the original network at $-x$,
the source is even, the data set is antipodally invariant, and the prior
is invariant. Draw $L$ independent networks from this one posterior.
Their ordinary raw features are then concatenated, with no even
projection, added input skip, or deterministic reflected feature copy.

The deterministic near-minimum-energy transfer certificate gives
comparators (4) for all $T=\binom r{2k}$ tasks. Their contamination is
linear in the input, so $d=r$. More explicitly, for deleted set $A$,
\[
 q_A(x)=\frac{\alpha}{N}\sum_i a_i\chi_A(z_i)
                               \relu(b_i+w_i\cdot x),
 \qquad
 \ell_A(x)=\frac{\alpha}{2N\sqrt r}
                       \sum_i a_i\|U_i\|\chi_A(z_i)z_i\cdot x.
 \tag{7}
\]
Here $z_i$ is the coordinate sign vector of $w_i$ and $\alpha$ is the
coefficient
\[
 \rho=\frac{\binom{r-2}{(r-2)/2}}{2^{r-1}\sqrt r},\qquad
 b_j=\rho(-1)^{(j-2)/2}
       \frac{\binom{(r-2)/2}{(j-2)/2}}{\binom{r-2}{j-2}},\qquad
 \alpha=\frac{b_r}{s b_{r-2k}}.
\]
Since $|A|$ is
even, (6) sends $\ell_A$ to $-\ell_A$. Zero weight coordinates have
posterior probability zero; an equivariant choice there is unnecessary.
Conditioning further on a common energy and population-hinge bound
preserves (6). Thus the balancing theorem applies conditionally on the
data and on every draw satisfying those bounds.

\section{Finite data and finite prior scale}
The earlier Gaussian transfer theorem supplies the following two facts.
First, at every data history, its stated width threshold and cold scale
give posterior energy failure probability at most $\beta^N$.
Second, with probability at least $1-\zeta_0$ over the $m$ original IID
inputs, every bounded-energy network fitting them has population hinge
loss at most $h_m$. The paired posterior fits those original inputs, so
the same uniform-convergence event applies without modification. Its
evidence lower bound used a full-margin witness and therefore applies
to the larger paired history as well.

For precision, its sufficient constants are as follows. Set
$D=r+1$, $p=r+2$, $E_*=s/\rho$, $E_0=(1+\varepsilon/4)E_*$,
and $E_{\max}=(1+\varepsilon)E_*$, with $0<\varepsilon\le1$. One may take
\[
 N_0=\left\lceil\frac{128E_0^2}{\varepsilon^2s^2}
                                    \log(4\cdot2^r)\right\rceil,
\]
\[
 a_{\rm box}=\min\left\{1,\frac{\varepsilon\sqrt{E_*}}{32\sqrt p},
 \frac{\varepsilon s}{32\sqrt D\{\sqrt{E_0}+(\sqrt{E_0}+1)\sqrt D\}}\right\},
 \qquad L_0=\log^+\frac{\sqrt{2\pi}}{2a_{\rm box}},
\]
\[
 \tau_0^2=\min\left\{1,
 \frac{\varepsilon E_*}
 {4\{pL_0+(p/2)\log(8/\varepsilon)+\log(1/\beta)\}}\right\}.
\]
\[
 h_m=\frac{8E_{\max}\sqrt D}{\sqrt m}
 +(s+E_{\max}\sqrt D)\sqrt{\frac{\log(1/\zeta_0)}{2m}}.
\]
Every $N\ge N_0$ and Gaussian standard deviation $0<\tau\le\tau_0$
obey the stated uniform posterior energy estimate. The deterministic
comparator constants in (4) can then be chosen as
\[
 \begin{split}
 u&=\sqrt{(s+E_{\max}\sqrt D)(\varepsilon s+2h_m)},\qquad
 v=\sqrt{2E_{\max}(\varepsilon s+h_m)/\rho},\\
 \Delta&=|\alpha|\{[1+(r-1)^k]v+(r-1)^ku\},\qquad
 B^2=2E_{\max}\alpha^2,\qquad R_0^2=2E_{\max}D.
 \end{split}
\]
These are the explicit inequalities behind the orders below, rather
than an unspecified limiting exchange of width and prior scale.

Consequently the ensemble good event has failure probability at most
\[
                         \zeta_0+L\beta^N+\zeta_1,     \tag{8}
\]
where $\zeta_1$ is the balancing failure in (3). On its complement,
(5) holds simultaneously for every nearby target. In particular there
is no condition $r/\binom r{2k}\le c\zeta\delta^2$.

For fixed $k,s$ and downstream accuracy $\delta$, choose the source
energy deficit $\varepsilon=c_{k,s}\delta^2r^{-4k-2}$ and
$h_m\le s\varepsilon$, with a sufficiently small constant. The
deterministic scales give
\[
 B^2=O_{k,s}(r^{2k+1}),\quad R_0^2=O_s(r^2),\quad
 \Delta\le\delta/8.
\]
Sufficient orders for the width of each bank and number of original
source inputs remain
\[
 N_0=O_{k,s}(r^{8k+7}\delta^{-4}),\qquad
 m=O_{k,s}(r^{8k+7}\delta^{-4}\log(2/\zeta_0)),        \tag{9}
\]
with $L=O(r+\log(T/\zeta_1))$. The total hidden width is $LN$.
Empirical hinge minimization over the radius-$2B$ ball of
$\mathcal H(\overline K)$, using
\[
 n=O_{k,s}(r^{2k+3}\delta^{-2}\log(2/\zeta_2)),        \tag{10}
\]
independent downstream examples, gives classification error at most
$\delta$ with probability at least
$1-\zeta_0-L\beta^N-\zeta_1-\zeta_2$, after adjusting the absolute
constants. Indeed the comparator's hinge loss is at most its $L^2$
error $\Delta$, because the target takes values in $\{-1,1\}$.
For separate downstream samples for all $T$ targets replace
$\zeta_2$ by $\zeta_2/T$.
If desired, choose $\beta$ so that $L\beta^{N_0}\le\zeta_0$;
the explicit prior-scale ceiling above then handles arbitrarily small
confidence levels without an additional width condition.

\begin{corollary}[Two raw banks and one posterior draw]\label{cor:antithetic}
On the paired source data above, draw one posterior network and append
its deterministic reflected bank. With the same $N_0,m,\tau_0$ and
comparator constants, all $T$ nearby tasks simultaneously have error
at most $\Delta$ and squared comparator norm at most $B^2$ in the
averaged two-bank kernel, except with probability
$\zeta_0+\beta^N$. Each of the two networks has the correct posterior
marginal. Their total hidden width is $2N$.
Hinge fitting in the radius-$B$ ball has the downstream sample order
(10), with total failure at most $\zeta_0+\beta^N+\zeta_2$ for a fixed
task; use $\zeta_2/T$ for separate samples on every task.
\end{corollary}
\begin{proof}
One good network supplies all decompositions (7). The targets are even,
the contamination is linear and odd, and the uniform cube measure is
invariant under $x\mapsto-x$. Apply Theorem~\ref{thm:antithetic}
without any additional random event. Reflection preserves the source
margin, energy, and population-hinge bounds. Equation (6) of the source
symmetry section gives the posterior marginal of the reflected network.
Only the one-draw energy failure and common data failure are needed.
The downstream argument is the preceding hinge comparator bound with
radius $B$ in place of $2B$.
\end{proof}

The antithetic corollary does not identify the concatenated pair with a
single width-$2N$ posterior draw: its components are deliberately
correlated. Its deterministic transfer conclusion also holds for a
non-Bayesian trained network satisfying the same energy and loss
hypotheses. In that application, reflection supplies the second bank
without requiring an invariant optimization distribution.

\paragraph{Equivariant optimization also supplies the symmetry.}
The balancing theorem is not specific to Bayesian training. On paired
even-source data, the empirical loss and Euclidean parameter penalty
are invariant under (6). Full-batch gradient descent with a fixed step
schedule therefore commutes with that transformation wherever the
gradient exists. The same is true with a fixed derivative convention
at ReLU zero and minibatches containing complete antipodal pairs.
Independent invariant initializations then produce independent,
conditionally symmetric contamination vectors. If such an optimization
procedure also attains the energy and population-loss hypotheses of
the deterministic comparator, the same ensemble conclusion follows.
This is a conditional implication; convergence to the required energy
level is not established here.

\paragraph{Scope.}
The independent-bank protocol uses a collection of independent
posterior draws and antipodally paired source observations to obtain
conditional symmetry. The separate antithetic protocol appends a
deterministic reflected bank to one draw. It does not establish the same simultaneous
guarantee for one unmodified posterior draw. The ordinary IID alternative
is to train each bank on an independent IID source data set; symmetry
then holds for the joint data--posterior law, at the cost of separate
source data. Exact independent posterior sampling remains an assumption.
The geometric and statistical arguments are conventional proofs, not
complete Lean formalizations.
\end{document}
